无限均匀双导线单位长度的电阻、电感、电导、电容分别为 \(R',L',G',C'\)（\(L',C'>0\)，\(R',G'\ge0\)）。令沿 \(+z\) 的电压、电流相量时间因子为 \(e^{i\omega t}\)，右端 \(z=0\) 接负载阻抗 \(Z_L\)，线沿 \(z\le0\) 延伸。考虑传播常数实部 \(\ge0\) 的分支，先对一般有损线求解，再取无损/弱损极限。
(1) 推导传输线相量形式微分方程、复传播常数 \(\gamma\) 与特性阻抗 \(Z_0\)。
(2) 求入射、反射振幅之比 \(\Gamma_L\) 及完全匹配条件，并说明有损情形与实阻抗的差异。
(3) 在无损情况下证明反射能流与透射负载功率满足守恒，列出开路、短路及匹配边界。
(4) 在 \(R'\ll\omega L'\)、\(G'\ll\omega C'\) 时一阶展开衰减常数及相位常数，并说明近似适用范围。
